\begin{quote}
\textbf{Maturity: {[}Partially Implemented{]}, and a design that
shifted.} This paper has changed more than any other in the series
between v2 and v3. The February 2026 draft proposed a specific
mechanism, an ERC-6551 NFT ``bridge-as-fundraise'' called \$SNAP, that
was \textbf{never deployed}. This revision keeps the durable thesis
(separate the instrument that raises liquidity from the gas token) and
replaces the unbuilt bridge with the money-path that Citrate actually
implemented: a wrapped SALT token with gasless authorization transfers,
a contracted market-maker allocation, and an on-chain SALT/USD price
feed. We say plainly what was proposed, what was built, and what is
still pending.
\end{quote}

\subsection{Abstract}\label{abstract}

Traditional blockchain fundraising creates a structural misalignment:
tokens sold to raise capital face immediate sell pressure, diluting the
network's economic foundation before it can build value. This paper's
durable claim is that the instrument used to bootstrap external
liquidity should be \emph{separate} from the native gas token, so that
raising and liquidity-provision do not dilute the utility asset. The
February 2026 draft realized that separation through \$SNAP, ERC-721
NFTs with ERC-6551 Token Bound Accounts serving simultaneously as
fundraising instruments, bridge nodes, and ownership positions. That NFT
bridge was specified and partially deployed to a testnet but never
carried to production, and it is not the mechanism Citrate runs. This
revision documents the implemented money-path: \textbf{WrappedSALT}, an
ERC-20 wrapper over native SALT supporting EIP-3009 gasless
authorization transfers; \textbf{MarketMakerAllocation}, a governable
vault that funds a contracted market-maker; and
\textbf{ComputePricingOracle}, which carries an on-chain SALT/USD price.
We present both, the design that shifted and the mechanism that exists,
and we keep the memetic-money framing as a philosophy rather than a
prediction. In keeping with Citrate Inc.'s public-communication policy,
this paper makes no claim of day-one cash earnings and treats SALT's
availability on a licensed exchange as pending.

\textbf{Keywords:} wrapped token, EIP-3009, gasless transfer, market
maker, price oracle, decentralized liquidity, memetic finance, token
economics, Citrate Network

\subsection{1. Introduction: the ICO
paradox}\label{introduction-the-ico-paradox}

Every network faces the same bootstrapping problem: it needs capital to
build infrastructure, but selling its native token to raise capital
creates sell pressure that undermines the token's economic foundation.
ICOs, IDOs, and venture rounds all involve selling or diluting the gas
token, so the network's economic base is weakened before it has
demonstrated value, and speculative holders exit precisely when the
network needs momentum.

The durable thesis of this paper is a structural response: separate the
instrument that bootstraps external liquidity from the gas token
entirely. Citrate's native gas token SALT (Paper I §5) should not be the
thing sold to raise or seed liquidity; a distinct mechanism should carry
SALT to external markets and provide the liquidity a two-sided market
needs, while SALT itself remains a utility asset. Everything else in
this paper is a question of \emph{which} mechanism realizes that
separation, and that is exactly where the design changed.

\subsection{2. The design that shifted: \$SNAP as an NFT
bridge-as-fundraise (February 2026, not
deployed)}\label{the-design-that-shifted-snap-as-an-nft-bridge-as-fundraise-february-2026-not-deployed}

\textbf{{[}Superseded, not deployed to production{]}} The February 2026
proposal realized the separation through \$SNAP: ERC-721 NFTs whose
ERC-6551 Token Bound Accounts held a dual-zone vault, a user-accessible
zone for the holder's fee revenue and a protocol-managed zone holding
locked ETH collateral that backed cross-chain transfers. Each NFT was to
be a bridge node; a per-wallet bonding curve discouraged whale
dominance; Chainlink oracle nodes operated by holders formed the
attestation layer; and the ETH committed at mint became bridge
liquidity, so that ``the fundraise is the bridge.'' A reference
implementation was deployed to an Ethereum testnet.

The design was intellectually appealing for reasons worth recording,
because they explain what the replacement had to preserve. The dual-zone
vault solved a real problem: it let a contributor retain ownership of
and withdraw their fee revenue (the user zone) while the collateral
backing transfers (the protocol zone) stayed locked behind a timelock
and governance, so holders could not drain the liquidity that made the
instrument useful. The per-wallet bonding curve, where a wallet's
successive mints grew progressively more expensive, was an attempt to
discourage whale dominance and reward early conviction without a central
allocator. And making each NFT a bridge node meant the act of investing
was the act of building capacity. These are genuine design virtues, and
any replacement had to keep the core of them: separate the roles of
``raise/seed liquidity'' and ``spend gas,'' and keep the liquidity base
stable.

We record the design here because honesty about the series' evolution
requires it, and because the thesis it embodied survives, but we are
explicit: \textbf{no SNAP/SALT lock-mint bridge contract exists on the
Citrate chain-40204 deployment.} The \texttt{CitrateBridgeNFT} /
\texttt{DualVault6551} / \texttt{BridgeRouter} contracts were never
carried to production, the AngelList SPV and Delaware C-Corp operating
structure described in the draft is not the vehicle for this work, and
cross-chain bridges, being among the most attacked primitives in the
space (Ronin \textasciitilde\$625M, Wormhole \textasciitilde\$325M,
Nomad \textasciitilde\$190M), are precisely the kind of surface a small
team should hesitate to ship without a very strong reason. The
per-wallet bonding curve also never solved Sybil resistance on its own,
a coordinated actor could mint once from many wallets at the base price,
so its whale-resistance depended on an off-chain identity layer the
draft acknowledged it lacked. Weighed together, the bridge's attack
surface and the instrument's unresolved Sybil problem argued for a
simpler mechanism. Section 3 documents the one that was built.

\subsection{3. What was actually built: the implemented
money-path}\label{what-was-actually-built-the-implemented-money-path}

The separation thesis is realized on chain 40204 by three deployed
contracts, none of which is a lock-mint bridge.

\textbf{3.1 WrappedSALT (EIP-3009 gasless transfer).}
\textbf{{[}Implemented, deployed \texttt{0xaa918302…}{]}}
\texttt{WrappedSALT.sol} is an ERC-20 wrapper over native SALT
implementing the EIP-3009 transfer-with-authorization flow:
\texttt{transferWithAuthorization}, a fee-bearing
\texttt{transferWithFeeAuthorization},
\texttt{receiveWithAuthorization}, \texttt{cancelAuthorization}, and
\texttt{authorizationState}, with a signature-malleability guard
(\texttt{\_recoverCanonical}). This is the ``portal'' in its concrete
form: a user signs an authorization off-chain, and a relayer or
counterparty submits it, so SALT can move to and be used by external
tooling without the holder needing native gas in hand. Where the
February design imagined an NFT-vault bridge, the implemented onramp is
a standard, auditable wrapped token with gasless authorization, the same
EIP-3009 pattern used by production stablecoins, plus native
\texttt{x402} payment-authorization precompiles (\texttt{0x0200},
\texttt{0x0201}) for the verification path.

\textbf{3.2 MarketMakerAllocation (contracted market maker).}
\textbf{{[}Implemented, deployed \texttt{0xa87fae5c…}{]}} Rather than
crowd-sourcing bridge liquidity from NFT holders, Citrate funds a
\emph{contracted} market-maker. \texttt{MarketMakerAllocation.sol} is a
\texttt{Governable} vault that accumulates native SALT (via a receive
fallback), and exposes \texttt{withdraw}/\texttt{withdrawAll},
\texttt{changeMarketMaker}, \texttt{changeAllocationRate}, and
\texttt{calculateAllocation}. The market-maker relationship, and its
allocation rate, are governance parameters, not a bonding curve. This is
the recast the series index anticipated: a contracted full-time
market-maker model in place of the NFT-node model.

\textbf{3.3 ComputePricingOracle (SALT/USD).} \textbf{{[}Implemented,
deployed \texttt{0xdcebd5ec…}{]}} There is no standalone
``SaltUsdOracle''; the SALT/USD reference lives in
\texttt{ComputePricingOracle.sol}, which carries
\texttt{saltPriceUsdCents} and a rate-limited \texttt{proposeSaltPrice}
update path (with a commit nonce and
\texttt{SaltPriceProposed}/\texttt{SaltPriceUpdated} events). This is
the price reference the allocation and compute-pricing logic read; it is
an operator-proposed feed with rate limiting, not a decentralized-oracle
network, and we describe it as such.

\textbf{3.4 The end-to-end money-path.} The three contracts compose into
a concrete flow that plays the role the February design assigned to the
bridge. A holder who wants to move SALT into external tooling wraps
native SALT into WrappedSALT and signs an EIP-3009 authorization; a
relayer or counterparty submits the authorization on-chain, so the
transfer settles without the holder paying native gas, and the
\texttt{x402} precompiles verify the authorization on the payment path.
On the market-making side, the treasury funds
\texttt{MarketMakerAllocation} with SALT at a governance-set allocation
rate, the contracted market-maker draws against that allocation to quote
two-sided liquidity on venues, and \texttt{ComputePricingOracle}
supplies the SALT/USD reference the allocation and compute-pricing logic
read. The separation the thesis demands is preserved: SALT reaches
external markets and a two-sided market is seeded without SALT being
sold in a fundraise, and the liquidity base is stable because it is a
governed allocation rather than a crowd of individually withdrawable
vaults. The difference from the February design is that the stability
now comes from a contracted relationship under governance control, not
from a timelocked NFT vault, and the whale-resistance question is moot
because there is no public mint to game. What the wrapped-token path
gives up relative to the NFT design is the ``every holder is an
operator'' ownership story; that ownership logic now lives in the
cooperative accounting of Paper VII rather than in bridge nodes.

\subsection{4. Security considerations}\label{security-considerations}

The shift in mechanism also shifts the threat model, mostly in a
favorable direction. \textbf{Bridge risk is largely avoided:} because
there is no lock-mint bridge holding pooled collateral, the catastrophic
bridge-exploit surface that dominated the February design (and the
industry's loss record) is simply not present. \textbf{The residual
risks are different and named.} The wrapped-token path inherits ERC-20
and EIP-3009 risks, replay and authorization-reuse, which the
implemented \texttt{cancelAuthorization}/\texttt{authorizationState} and
the canonical-signature recovery are there to bound, and it should still
be third-party audited before any mainnet reliance. The market-maker
path introduces a \emph{trust} dependency:
\texttt{MarketMakerAllocation} funds a contracted counterparty, and the
governance controls (\texttt{changeMarketMaker},
\texttt{changeAllocationRate}) are the only checks on that relationship,
so their timelock and quorum matter. The price feed is operator-proposed
and rate-limited, which is honest to state: it is not manipulation-proof
the way a deep decentralized oracle would be, and downstream logic
should treat it accordingly. Two specifics are worth stating because
they are the failure modes a reviewer will look for first. On the
wrapped-token side, EIP-3009 authorizations are single-use nonced
messages: each carries a unique nonce whose consumption is recorded in
\texttt{authorizationState}, so a submitted authorization cannot be
replayed, and a holder who signed one they no longer want executed can
burn it via \texttt{cancelAuthorization} before a relayer submits it.
The malleability guard matters here too, because without
canonical-signature recovery an attacker could reshape a valid signature
into a second distinct-looking authorization;
\texttt{\_recoverCanonical} closes that. On the market-maker side, the
honest statement is that \texttt{MarketMakerAllocation} concentrates
trust in a contracted counterparty: the vault can be drained up to its
allocation by whoever holds the market-maker role, so the safety of the
arrangement reduces to the safety of \texttt{changeMarketMaker} and
\texttt{changeAllocationRate}, which are the governance levers that
appoint and bound that counterparty. Those levers should carry the same
timelock and quorum as any treasury action (Paper VIII), and their event
log is the audit trail a holder uses to verify the market-maker has not
been changed without notice. As with every contract in the series, none
of this substitutes for a professional audit, which remains a
precondition for mainnet.

\subsection{5. The memetic thesis}\label{the-memetic-thesis}

The name ``Memetic Money Portal'' is deliberate, and survives the
mechanism change. In traditional finance, a monetary instrument's value
derives from institutional backing; in cryptocurrency, much of it
derives from narrative propagation. The meme-coin phenomenon showed that
cultural resonance can bootstrap real economic networks, though meme
coins typically lack utility beyond speculation. The portal's aspiration
is to couple narrative and utility: the instrument that spreads
culturally is also the instrument that does useful work, moving value
and seeding liquidity. \textbf{Honest boundaries:} this depends on
achieving cultural resonance, which is not an engineering variable and
cannot be specified or guaranteed. We present the memetic thesis as a
design philosophy, not a prediction, and we note it is even more clearly
a philosophy now that the mechanism is a wrapped token and a
market-maker rather than a collectible NFT.

\subsection{6. Compliance and honest
boundaries}\label{compliance-and-honest-boundaries}

This paper is written under Citrate Inc.'s public-communication policy,
and two constraints apply directly. First, \textbf{no day-one cash
earnings} are claimed from holding SALT, providing liquidity, or
operating any part of this money-path; SALT is a utility and staking
asset, cash convertibility is gated on mainnet and on audit, and any
yield is a staked-grant construct, not a cash promise. Second,
\textbf{SALT's availability on a licensed exchange is pending}, so
nothing here should be read as a listing commitment or a solicitation.
The February draft's fundraise targets, hard caps, and SPV structure are
not carried forward; the implemented money-path is infrastructure, not
an offering.

\subsection{7. Relationship to the Gradient Papers
Series}\label{relationship-to-the-gradient-papers-series}

Paper I defines SALT tokenomics (1B supply, distribution, staking
minimums); the Memetic Money Portal provides the mechanism by which SALT
reaches external markets without being sold in a fundraise. Paper VII
(Mozi Cooperative) provides the economic philosophy of
contribution-proportional ownership; the implemented money-path is
thinner than the NFT-holder model the draft imagined, and the
cooperative's ownership logic now lives primarily in
\texttt{ContributionAccounting} (Paper VII), not in bridge NFTs. Paper
VIII (BR1J Constitution) supplies the governance that controls the
market-maker and price-feed parameters through
\texttt{TreasuryGovernor}. Paper IX (Medusa Paradigm) offers the
symbiotic framing, a host and its guest each benefiting, which now maps
to the network and its contracted market-maker rather than to a
coral-and-algae vault.

\subsection{8. Conclusion}\label{conclusion}

The Memetic Money Portal's durable idea, separate the instrument that
bootstraps liquidity from the gas token, holds. Its \emph{mechanism}
changed: the ERC-6551 NFT bridge-as-fundraise of February 2026 was never
deployed, and the money-path Citrate actually runs is a wrapped SALT
token with EIP-3009 gasless transfers, a governable market-maker
allocation, and an on-chain SALT/USD feed, all deployed on chain 40204.
This is a narrower, more auditable, and less dangerous design than a
pooled-collateral bridge, and documenting the change honestly is the
point of this revision. What remains before mainnet reliance is the same
short list that gates the rest of the series: a professional audit, real
SALT liquidity, and, for anything touching users' funds, the compliance
posture Section 6 describes.

\subsection{Acknowledgments}\label{acknowledgments}

Drafting and literature triage were assisted by AI systems; all
mechanical claims were verified by the authors against the referenced
source files and deployed contracts on chain 40204. This work received
no external funding.

\subsection{References}\label{references}

{[}1{]} ERC-6551: Non-fungible Token Bound Accounts. Ethereum
Improvement Proposal (Windle, Giang, et al., 2023). \emph{(Cited for the
superseded design of Section 2.)} {[}2{]} ERC-3009: Transfer with
Authorization. Ethereum Improvement Proposal. {[}3{]} ERC-20: Token
Standard. Ethereum Improvement Proposal (Vogelsteller, Buterin, 2015).
{[}4{]} Klosowski, L., Mendenhall, L. (2026). Citrate: Protocol
Specification for an AI-Native BlockDAG Network. \emph{The Gradient
Papers No.~I} (this series). {[}5{]} Klosowski, L., Mendenhall, L.
(2026). The Mozi Cooperative. \emph{The Gradient Papers No.~VII} (this
series). {[}6{]} Klosowski, L., Mendenhall, L. (2026). The BR1J
Constitution. \emph{The Gradient Papers No.~VIII} (this series). {[}7{]}
Klosowski, L., Mendenhall, L. (2026). The Medusa Paradigm. \emph{The
Gradient Papers No.~IX} (this series). {[}8{]} Buterin, V. (2014).
Ethereum: a next-generation smart contract and decentralized application
platform. {[}9{]} Ronin Network (2022); Wormhole (2022); Nomad (2022).
Bridge incident post-mortems. \emph{(Cited for the bridge-risk record
that motivated not shipping a lock-mint bridge.)}

\subsection{Appendix A: Cross-Paper Parameter Consistency (reconciled
against code, Aug
2026)}\label{appendix-a-cross-paper-parameter-consistency-reconciled-against-code-aug-2026}

{\def\LTcaptype{none} % do not increment counter
\begin{longtable}[]{@{}
  >{\raggedright\arraybackslash}p{(\linewidth - 4\tabcolsep) * \real{0.4231}}
  >{\raggedright\arraybackslash}p{(\linewidth - 4\tabcolsep) * \real{0.2692}}
  >{\raggedright\arraybackslash}p{(\linewidth - 4\tabcolsep) * \real{0.3077}}@{}}
\toprule\noalign{}
\begin{minipage}[b]{\linewidth}\raggedright
Parameter
\end{minipage} & \begin{minipage}[b]{\linewidth}\raggedright
Value
\end{minipage} & \begin{minipage}[b]{\linewidth}\raggedright
Source
\end{minipage} \\
\midrule\noalign{}
\endhead
\bottomrule\noalign{}
\endlastfoot
SALT total supply & 1,000,000,000 & Paper I §5 \\
BFT threshold & 67/100 (2/3) & Paper I; \texttt{checkpoint.rs:88} \\
WrappedSALT (EIP-3009) & deployed & \texttt{0xaa918302…};
\texttt{WrappedSALT.sol} \\
MarketMakerAllocation & deployed & \texttt{0xa87fae5c…};
\texttt{MarketMakerAllocation.sol} \\
SALT/USD feed (ComputePricingOracle) & deployed & \texttt{0xdcebd5ec…};
\texttt{ComputePricingOracle.sol} \\
x402 payment-auth precompiles & \texttt{0x0200}--\texttt{0x0201} &
\texttt{precompiles/x402.rs} \\
SNAP/SALT lock-mint bridge & does not exist & superseded design (Section
2) \\
\end{longtable}
}

\begin{center}\rule{0.5\linewidth}{0.5pt}\end{center}

\emph{This paper is part of the Gradient Papers series, published by
Citrate Inc.} \emph{Correspondence: Larry@citrate.ai}
